% This is a variation of the NeurIPS'24 format
\documentclass{article}
\usepackage[final]{adrl}

\usepackage[utf8]{inputenc} % allow utf-8 input
\usepackage[T1]{fontenc}    % use 8-bit T1 fonts
\usepackage{hyperref}       % hyperlinks
\usepackage{url}            % simple URL typesetting
\usepackage{booktabs}       % professional-quality tables
\usepackage{amsfonts}       % blackboard math symbols
\usepackage{nicefrac}       % compact symbols for 1/2, etc.
\usepackage{microtype}      % microtypography
\usepackage{xcolor}
\usepackage{algorithm}
\usepackage{algpseudocode}

\title{When and Why Does Hindsight Experience Replay\\
Improve Soft Actor-Critic under Sparse Rewards?}
\author{Mustafa Asfari \\ Mohammad Alkhozami \\
\url{github.com/mustafa99705/adrl-sac-her-fetchreach}}

\begin{document}

\maketitle

\section{Motivation}
In real robotic tasks, the most natural reward specification is often binary:
the task was either accomplished or it was not. Dense reward functions can
provide more feedback, but they usually have to be hand-engineered, require
domain knowledge, and may bias the learned behavior in unintended ways.
However, standard off-policy actor-critic methods such as Soft Actor-Critic
(SAC) can struggle under sparse rewards, because the replay buffer may contain
many uninformative failed transitions before the agent observes its first
success. While Hindsight Experience Replay (HER) has been shown to improve
learning in sparse-reward settings, it is less clear under which task
conditions its benefit is most significant, and through which mechanism it
acts. This project therefore investigates not only \emph{whether} HER improves
SAC, but \emph{when} its contribution becomes meaningful as task difficulty
increases, and \emph{why} it helps, by analyzing the training dynamics it
induces.

\section{Related Topics}
This project connects three main topics from the lecture and the RL
literature: off-policy maximum-entropy reinforcement learning,
goal-conditioned reinforcement learning, and reward design under sparse
feedback. SAC~\cite{haarnoja18} is the off-policy actor-critic algorithm used
as the learning backbone, while HER~\cite{andrychowicz17} provides the
goal-relabeling mechanism that allows failed trajectories to be reused with
achieved goals. The project also relates to Universal Value Function
Approximators~\cite{schaul15} and multi-goal RL, because both the policy and
the value function must condition on the desired goal. The Fetch environments
and their sparse/dense reward variants are based on the multi-goal robotics
benchmark introduced by Plappert et al.~\cite{plappert18} and are available
through Gymnasium-Robotics.

\section{Idea}
The idea of this project is to investigate under which task conditions HER
provides the greatest benefit for SAC in goal-conditioned continuous control,
and to shed light on the mechanism behind that benefit. Rather than focusing
only on whether HER improves performance, we compare its effect across
robotic manipulation tasks of different difficulty. We will evaluate four
experimental conditions --- dense SAC, dense SAC+HER, sparse SAC, and sparse
SAC+HER --- first on \texttt{FetchReach-v3}, a relatively simple reaching
task, and then on \texttt{FetchPush-v3}, where the robot must push an object
to a target location and sparse-reward learning is known to be substantially
harder~\cite{andrychowicz17}. This controlled setup allows us to study
whether the benefit of HER grows with task difficulty while keeping the
learning algorithm unchanged. Our hypotheses: (i) on the simple reaching
task, HER will provide only a modest speedup, since random exploration
occasionally reaches the goal on its own; (ii) on the harder pushing task,
HER will make the difference between learning and not learning at all;
(iii) HER's effect will be visible in the training dynamics, e.g.\ in the
evolution of SAC's entropy coefficient and in the distribution of relabeled
goals, which is expected to form an implicit curriculum.

\begin{algorithm}[H]
    \caption{SAC with optional Hindsight Experience Replay (\emph{future} strategy)}
    \label{alg:code}
    \begin{algorithmic}
        \Require environment $e$, SAC algorithm $A$, reward function $r$,
                 HER flag $h$, relabels per transition $k$
        \Ensure trained goal-conditioned policy $\pi(a \mid s, g)$
        \State Initialize replay buffer $R$, SAC actor and critic networks
        \While{training budget not reached}
            \State Sample goal $g$ and initial state $s_0$ from $e$
            \State Roll out one episode $(s_0, a_0, \dots, s_T)$ with $\pi(\cdot \mid s, g)$
            \State Store $(s_t, a_t, r(s_t, a_t, g), s_{t+1}, g)$ in $R$ for all $t$
            \If{$h$ is true}
                \For{each $t$ and each of $k$ relabels}
                    \State sample $\tau \sim \mathcal{U}(t+1, T)$;\;
                           $g' \gets$ achieved goal in $s_\tau$
                    \State store $(s_t, a_t, r(s_t, a_t, g'), s_{t+1}, g')$ in $R$
                           \Comment{often a success: $r = 0$}
                \EndFor
            \EndIf
            \State Update SAC actor and critics on minibatches from $R$
        \EndWhile
        \State \Return $\pi$
    \end{algorithmic}
\end{algorithm}

\section{Experiments}
\paragraph{Environments \& Metrics}
We will use two goal-conditioned continuous-control environments from
Gymnasium-Robotics: \texttt{FetchReach-v3} as a simple reference task and
\texttt{FetchPush-v3} as the main, more challenging benchmark, since the
robot must manipulate an object toward a desired goal. For both tasks we
consider the sparse reward variant, with the corresponding dense-reward
variants as reference baselines. The experiments will use SAC from
Stable-Baselines3~\cite{raffin21sb3} with a \texttt{MultiInputPolicy} for the
dictionary observations. Every 2{,}000 training steps we will evaluate the
deterministic policy on 20 held-out episodes. The primary metric is the
\textbf{success rate}, using the success criterion defined by each
environment (end effector, respectively object, within 5\,cm of the goal at
episode end). From the success-rate curves we will report the median number
of environment steps to reach 90\% success, the area under the success curve
as a measure of sample efficiency, and the final success rate. Episode
return and training-time success rates will be logged as secondary metrics.

\paragraph{Experimental Scope}
The study covers four conditions: dense SAC, dense SAC+HER, sparse SAC, and
sparse SAC+HER. On \texttt{FetchReach-v3} we will run all four conditions
with 10 random seeds each. On \texttt{FetchPush-v3}, which is computationally
heavier, we will focus on the two central conditions --- sparse SAC and
sparse SAC+HER --- again with 10 seeds. We deliberately avoid tuning
hyperparameters separately per condition, since per-condition tuning would
confound the effect of the reward signal and of HER with the effect of the
tuning itself. All runs therefore share one fixed configuration adopted from
the tuned settings of the RL Baselines3 Zoo for Fetch-style
tasks~\cite{raffin20zoo} ($\gamma = 0.95$, learning rate $10^{-3}$, batch
size 256, replay buffer size $10^6$; all remaining settings are
Stable-Baselines3 defaults~\cite{raffin21sb3}). Short pilot runs will verify
that this configuration transfers to our setup and will calibrate the
training budget per environment before the main grid is launched.

As a required ablation, we will vary the HER relabeling ratio,
$n_{\text{sampled\_goal}} \in \{1, 4, 8\}$, on \texttt{FetchPush-v3} with 5
seeds per setting, to test whether more hindsight relabeling helps
monotonically or saturates. To address the question of \emph{why} HER helps,
we will additionally track the evolution of SAC's automatically tuned entropy
coefficient $\alpha$ across conditions, and log the distribution of relabeled
goal distances over training as an indicator of the implicit curriculum
induced by HER.

\paragraph{Estimated Computational Load}
We will first perform short pilot runs to measure wall-clock training speed
and determine an appropriate step budget per environment; based on published
results we expect FetchReach to require on the order of $10^5$ environment
steps per run (roughly 20--60 minutes each) and FetchPush substantially more
(on the order of $10^6$ steps, i.e.\ a few hours per run). In total we
expect the complete study --- core grid plus ablation --- to require several
tens of GPU-hours. The experiments will be executed on our own available
hardware using Stable-Baselines3, with modest per-run requirements
(one GPU, 4 CPU cores, 16\,GB RAM, and a few gigabytes of storage for logs
and checkpoints); the runs are independent, so they parallelize trivially
across machines if needed.

\section{Timeline}
\paragraph{Research and Initial Implementation (week 1)}
We will review HER, goal-conditioned RL, and sparse-reward learning, and
study the FetchReach and FetchPush environment definitions (observation and
goal spaces, reward semantics, success criteria) as well as the relevant SAC
and HER implementation details in Stable-Baselines3. In parallel we will fix
the software versions and begin implementing the common training pipeline
with seeding, logging, evaluation, and checkpointing. Short smoke tests will
verify the environment setup and reward computation.

\paragraph{Implementation and Experiments (week 2)}
We will complete the training pipeline for all four conditions, run pilot
experiments to check stability and calibrate step budgets, and then launch
the main grid on \texttt{FetchReach-v3} and \texttt{FetchPush-v3}. After
validating the core comparison, we will run the HER relabeling-ratio
ablation. Throughout, we will log entropy-coefficient trajectories and
relabeled-goal statistics for the mechanism analysis.

\paragraph{Analysis and Reporting (week 3)}
We will aggregate results across seeds, inspect per-seed learning curves for
instability or failed runs, and compute success-rate curves,
steps-to-90\%-success, area under the success curve, and final success rate.
The analysis will focus on how the benefit of HER changes with task
difficulty and relabeling ratio, and on the mechanism indicators (entropy
coefficient evolution, relabeled-goal distributions). Finally, we will
prepare the figures, write the report, document the experimental setup, and
provide reproducibility instructions.

\begin{thebibliography}{9}
\bibitem{haarnoja18}
T.~Haarnoja, A.~Zhou, P.~Abbeel, S.~Levine.
\newblock Soft Actor-Critic: Off-Policy Maximum Entropy Deep Reinforcement
Learning with a Stochastic Actor.
\newblock \emph{ICML}, 2018.

\bibitem{andrychowicz17}
M.~Andrychowicz, F.~Wolski, A.~Ray, J.~Schneider, R.~Fong, P.~Welinder,
B.~McGrew, J.~Tobin, P.~Abbeel, W.~Zaremba.
\newblock Hindsight Experience Replay.
\newblock \emph{NeurIPS}, 2017.

\bibitem{plappert18}
M.~Plappert, M.~Andrychowicz, A.~Ray, B.~McGrew, B.~Baker, G.~Powell,
J.~Schneider, J.~Tobin, M.~Chociej, P.~Welinder, V.~Kumar, W.~Zaremba.
\newblock Multi-Goal Reinforcement Learning: Challenging Robotics Environments
and Request for Research.
\newblock \emph{arXiv:1802.09464}, 2018.

\bibitem{schaul15}
T.~Schaul, D.~Horgan, K.~Gregor, D.~Silver.
\newblock Universal Value Function Approximators.
\newblock \emph{ICML}, 2015.

\bibitem{raffin20zoo}
A.~Raffin.
\newblock RL Baselines3 Zoo: A Training Framework for Stable-Baselines3
Reinforcement Learning Agents.
\newblock \url{https://github.com/DLR-RM/rl-baselines3-zoo}, 2020.

\bibitem{raffin21sb3}
A.~Raffin, A.~Hill, A.~Gleave, A.~Kanervisto, M.~Ernestus, N.~Dormann.
\newblock Stable-Baselines3: Reliable Reinforcement Learning Implementations.
\newblock \emph{JMLR}, 22(268), 2021.
\end{thebibliography}

\end{document}
